\documentclass{leveldoc}
\DocUses{foundations,deflection_plates,tune_the_lens}
\DocFingerprint{66e88abbeccfddb9}

\begin{document}

\section{What the level asks}\label{sec:asks}

\leveldescription

Three positive ions are launched from the left edge, one straight along $+x$
and two at small angles to it. Three metal apertures stand on the launch
line. The outer apertures are grounded. The middle aperture is held at a
fixed potential $\val{vMid}$. That lens is the one in Tune the lens, at a
lower potential: it slows a positive ion and turns the angled rays back
toward the launch line, and the focus falls short of the detector. The player
adds charges downstream of the lens. The reference places two.

Each ion must enter the detector. No shot has an acceptance window, a kinetic
window or a gate. A flight ends when the ion touches a plate or a charge,
leaves the world, or enters the detector (\xsecref{foundations}{sec:board}).
A plate's obstacle is its box plus a contact shell of thickness
$\val{contact}$. A charge's obstacle has radius $\val{Rcharge}$. Each ion's
radius is $0$.

With no charges, the on-axis ray still enters the detector and the angled
rays miss it by a fraction of a cell. This document uses the deflection
plates' rectangle collocation, adds the Coulomb field of the placed charges,
and integrates the three flights by the classical Runge--Kutta method.

\section{Setup}\label{sec:setup}

\setuptable

The speed of light is $c = \val{cLight}$. Each ion has charge $q = \val{qPart}$,
mass $m = \val{mPart}$ and radius $0$, and each is launched with kinetic
energy $T_0 = \val{T0}$ from $(\val{xLaunch}, \val{yLaunch})$. Then
$\gamma_0 = \val{gamma0}$, $\abs{\vect{p}_0} = \val{p0}$ and
$v_0 = \val{v0} = \val{vOverC}\,c$ (\xeqref{foundations}{eq:launch}). The
three directions make angles $\val{aLow}$, $\val{aAxis}$ and $\val{aHigh}$
degrees with $+x$.

The apertures stand at $x = \val{xIn}$, $x = \val{xMid}$ and $x = \val{xOut}$,
with the same boxes as in Tune the lens. The opening between the inner rims
is $\val{opening}$, from $\val{yRimLo}$ to $\val{yRimHi}$. The middle halves
are at $V = \val{vMid}$. The four outer halves are at $\val{vOuter}$.

The reference charges are $q = \val{qC0}$ at $(\val{xC0}, \val{yC0})$ and
$q = \val{qC1}$ at $(\val{xC1}, \val{yC1})$, both below the launch line and
downstream of the lens. The charge magnitudes on offer run up to $\val{qMag}$.

The detector, the same one for every shot, is $\val{xDetLo} \le x \le \val{xDetHi}$,
$\val{yDetLo} \le y \le \val{yDetHi}$. The world is $-1 \le x \le \val{xWorld}$,
$-1 \le y \le \val{yWorld}$. The reported time step is $h = \val{H}$. The mesh
is the deflection plates' mesh: rectangles of edge $\val{panelFine}$,
$\val{nPanels}$ of them on the upper half.

\begin{figure}[ht]
\centering
\begin{tikzpicture}
\begin{axis}[axis equal image, width=0.92\linewidth,
  xmin=-1.5, xmax=32.5, ymin=-2, ymax=22.5,
  xlabel={$x$}, ylabel={$y$}, tick label style={font=\small}]
  \draw[black!40, dashed] (-1,-1) rectangle (\rawval{xWorld},\rawval{yWorld});
  \draw[black!25] (0,0) rectangle (30,20);
  \fill[green!50!black, fill opacity=0.18]
    (\rawval{xDetLo},\rawval{yDetLo}) rectangle (\rawval{xDetHi},\rawval{yDetHi});
  \draw[green!40!black]
    (\rawval{xDetLo},\rawval{yDetLo}) rectangle (\rawval{xDetHi},\rawval{yDetHi});
  \fill[black!20] (\rawval{xP0a},\rawval{yP0a}) rectangle (\rawval{xP0b},\rawval{yP0b});
  \fill[black!20] (\rawval{xP1a},\rawval{yP1a}) rectangle (\rawval{xP1b},\rawval{yP1b});
  \fill[black!20] (\rawval{xP4a},\rawval{yP4a}) rectangle (\rawval{xP4b},\rawval{yP4b});
  \fill[black!20] (\rawval{xP5a},\rawval{yP5a}) rectangle (\rawval{xP5b},\rawval{yP5b});
  \fill[black!20] (\rawval{xP2a},\rawval{yP2a}) rectangle (\rawval{xP2b},\rawval{yP2b});
  \draw[orange!80!black]
    (\rawval{xP2a},\rawval{yP2a}) rectangle (\rawval{xP2b},\rawval{yP2b});
  \fill[black!20] (\rawval{xP3a},\rawval{yP3a}) rectangle (\rawval{xP3b},\rawval{yP3b});
  \draw[orange!80!black]
    (\rawval{xP3a},\rawval{yP3a}) rectangle (\rawval{xP3b},\rawval{yP3b});
  \fill[violet!70!black] (\rawval{xC0},\rawval{yC0}) circle[radius=\rawval{Rcharge}];
  \fill[red!70!black] (\rawval{xC1},\rawval{yC1}) circle[radius=\rawval{Rcharge}];
  \addplot[orange!80!black] table {pathofflow.dat};
  \addplot[orange!80!black] table {pathoffhigh.dat};
  \addplot[orange!80!black] table {pathoffaxis.dat};
  \addplot[blue!60!black] table {pathlow.dat};
  \addplot[very thick, blue!60!black] table {pathaxis.dat};
  \addplot[blue!60!black] table {pathhigh.dat};
  \fill (\rawval{xLaunch},\rawval{yLaunch}) circle[radius=0.08];
\end{axis}
\end{tikzpicture}
\caption{The six aperture plates, the two charges and the six flights. The
outlined plates are the middle aperture, at the fixed potential. The violet
disc is the negative charge and the red disc is the positive charge. The
green box is the detector. Blue paths are the reference: the downward ray
meets the bottom edge, the on-axis ray meets the left edge, and the upward
ray meets the top edge at its left corner. Orange paths are the same launches
with the charges omitted. The on-axis orange path meets the left edge. The
other two stay clear of the detector and leave through the right edge of the
world, one just below the detector and one just above it.}
\label{fig:board}
\end{figure}

\section{Principles and model}\label{sec:model}

\begin{principles}
  \item The relativistic equation of motion \xeqref{foundations}{eq:eom}, the
    launch \xeqref{foundations}{eq:launch}, and energy in a static field
    \xeqref{foundations}{eq:energy}. There is no magnetic field, and the motion
    stays in the plane (\xsecref{foundations}{sec:plane}).
  \item A fixed potential on an electrode, and a grounded electrode at
    potential $0$ (\physref{2.7}). The surface charge is the deflection
    plates' collocation \xeqref{deflection_plates}{eq:collocation}, with the
    rectangle integral \xeqref{deflection_plates}{eq:rectangle}. A placed
    charge adds its Coulomb potential to the right-hand side and its Coulomb
    field to the flight.
  \item With no placed charge the lens is the symmetric aperture lens of
    Tune the lens (\xeqref{tune_the_lens}{eq:axis},
    \xeqref{tune_the_lens}{eq:slow}). The placed charges sit off that line,
    so the reference flights are not the symmetric ones.
  \item The board's rules (\xsecref{foundations}{sec:board}) and the classical
    Runge--Kutta method with cubic Hermite event location
    (\xsecref{foundations}{sec:numerics}).
  \item The radiated power \xeqref{foundations}{eq:larmor}, summed by Simpson's
    rule \xeqref{foundations}{eq:simpson}. It is not put back into the force.
\end{principles}

Off the metal the field is the field of the surface charge, of its mirror,
and of the placed charges. The state $(\vect{x}, \vect{p})$ obeys
\xeqref{foundations}{eq:eom} with $\vect{B} = \vect{0}$, from the launch
until the first ending event.

\begin{assumption}{The potential is static}\label{as:static}
The middle aperture is held at $\val{vMid}$ with no supply and no ramp. The
launch speed is $\val{vOverC}\,c$. The game uses the same static bias
(\physref{2.7}).
\end{assumption}

\begin{assumption}{No radiation reaction}\label{as:radiation}
Each ion radiates, and that energy is not taken from the flight. The largest
of the three reference integrals is a fraction $\val{radFracMax}$ of $T_0$.
\end{assumption}

\begin{assumption}{No image of the ion}\label{as:image}
The ion's own surface charge is left out, as the deflection plates leave it
out and as the game leaves it out (\physref{2.7}). The neglected force is
bounded by $q^2/d^2$, with the distance $d$ floored at $10^{-3}$, relative to
$\max(\abs{q\vect{E}}, T_0)$. The reference samples reach $\val{refImage}$,
under the engine's limit $10^{-10}$. With the charges omitted the samples
reach $\val{offImage}$. The samples are the step ends and the ending point.
\end{assumption}

\begin{assumption}{Equal rectangles}\label{as:panels}
The mesh is the one stated for the deflection plates: equal rectangles of edge
$\val{panelFine}$, not the engine's graded triangles. Replacing the edge by
$\val{panelCoarse}$ moves an end state by $\val{errMesh}$. Both placements
carry electrode charge, so the estimate uses every ray.
\end{assumption}

\begin{assumption}{The fixed step}\label{as:step}
The reported step is $h = \val{H}$. Doubling it estimates an end state by
$\val{errStep}$. Ten times the panel-size estimate $\val{errMesh}$ lies above
$10^{-9}$, and it is larger than the step estimate $\val{errStep}$, so the
tolerance on a state is that product, $\val{tolFlight}$. Between samples the
ion moves at most $\val{stepReach}$. The smallest sampled clearance is
$\val{refDmin}$ on the reference flights and $\val{offDmin}$ with the charges
omitted.
\end{assumption}

\section{Derivation}\label{sec:derivation}

\subsection{The lens without the charges}\label{sec:lens}

The middle halves are at $\val{vMid}$ and the outer halves are at
$\val{vOuter}$. With no placed charge the plates are symmetric about
$y = \val{yLaunch}$, so \xeqref{tune_the_lens}{eq:axis} applies. The panel
solution leaves a leftover $\val{eAxis}$ at the middle of the opening. An ion
launched on the line stays on it. \xeqref{tune_the_lens}{eq:slow} then gives
the kinetic energy. At the launch, $\phi = \val{phiLaunch}$. In the middle
opening, $\phi = \val{phiMid}$, so $T = \val{Tmid}$. The ion gives up a
fraction $\val{slowFrac}$ of $T_0$. On the launch line the potential is
$\val{phiIn}$ at the entrance aperture and $\val{phiOut}$ at the exit
aperture.

Above the line, at $(\val{xMid}, \val{yAbove})$, the panel field is
$\val{eAbove}$, toward the line. Below it, at $(\val{xMid}, \val{yBelow})$,
the field is $\val{eBelow}$, also toward the line. A positive ion is pushed
toward the launch line. At this potential the angled rays still miss the
detector, as the flights in \S\ref{sec:others} show.

\subsection{The charges and the panels}\label{sec:charges}

The reference adds the charges of \S\ref{sec:setup}. Each contributes
$\phi = q/r$ and $\vect{E} = q\,\hat{\vect{r}}/r^2$, with $k = 1$. At the
middle of the opening the transverse field, now including those charges, is
$\val{ePlaced}$. The collocation is
\xeqref{deflection_plates}{eq:collocation}: the right-hand side is the
electrode potential minus the charges' potential. The reported mesh has
$\val{nPanels}$ upper rectangles. The residual at the centroids, with the
charges included, is $\val{solveRes}$. On the large faces, off the centroids,
the potential differs from the electrode potential by at most $\val{surfRes}$.

With the charges omitted the middle aperture carries $\val{qMid}$ and the
grounded halves carry $\val{qOuter}$. The state is advanced by
\xeqref{foundations}{eq:rk4} with step $h = \val{H}$. Events are located as
on the deflection plates: thirty-two Hermite samples on the step, then sixty
bisections, obstacles before the world, and the world before the detector.
A face coordinate is assigned when the residual is below $10^{-6}$. The
step-doubling estimate divides the change by $15$; the panel-size estimate
is the change itself when the edge doubles.

\section{The reference solution by hand}\label{sec:hand}

\begin{procedure}
  \item\label{st:launch} \emph{The launches} \xeqref{foundations}{eq:launch}.
    $\gamma_0 = \val{gamma0}$ and $\abs{\vect{p}_0} = \val{p0}$, so
    $v_0 = \val{v0}$. The three angles are $\val{aLow}$, $\val{aAxis}$ and
    $\val{aHigh}$ degrees.
  \item\label{st:lens} \emph{The lens.} Hold both halves at $x = \val{xMid}$
    at $\val{vMid}$ and the other four electrodes at $\val{vOuter}$. The
    opening is $\val{opening}$. On the launch line, with no charges,
    \xeqref{tune_the_lens}{eq:slow} gives $T = \val{Tmid}$ in the middle
    opening.
  \item\label{st:charges} \emph{The charges.} Place $\val{qC0}$ at
    $(\val{xC0}, \val{yC0})$ and $\val{qC1}$ at $(\val{xC1}, \val{yC1})$.
  \item\label{st:panels} \emph{The panels.} Solve
    \xeqref{deflection_plates}{eq:collocation} on the mesh of
    \S\ref{sec:charges}, with the charges' potential on the right-hand side.
    The centroid residual is $\val{solveRes}$. The off-centroid discrepancy
    on the large faces is $\val{surfRes}$.
  \item\label{st:march} \emph{The marches.} Advance
    \xeqref{deflection_plates}{eq:flight-energy} and
    \xeqref{foundations}{eq:eom} by \xeqref{foundations}{eq:rk4} with step
    $\val{H}$, once for each launch direction. The on-axis state every five
    time units:
    \begin{center}\small\val{states}\end{center}
  \item\label{st:end} \emph{The arrivals.} The downward ray meets the
    \val{LowFace} face of the detector at $(\val{LowX}, \val{LowY})$, at
    $t = \val{LowT}$, with momentum $(\val{LowPx}, \val{LowPy})$, at an angle
    $\val{LowAng}$ degrees from $+x$. The ordinate is the detector's bottom
    edge. The on-axis ray meets the \val{AxisFace} face at
    $(\val{AxisX}, \val{AxisY})$, at $t = \val{AxisT}$, with momentum
    $(\val{AxisPx}, \val{AxisPy})$. The abscissa is the detector's left edge.
    The upward ray meets the \val{HighFace} face at
    $(\val{HighX}, \val{HighY})$, at $t = \val{HighT}$, with momentum
    $(\val{HighPx}, \val{HighPy})$, at an angle $\val{HighAng}$ degrees from
    $+x$. The ordinate is the detector's top edge, at its left corner. The
    runs take $\val{LowSteps}$, $\val{AxisSteps}$ and $\val{HighSteps}$ steps.
  \item\label{st:check} \emph{The checks.} Along the three flights, $W$ drifts
    by at most $\val{LowWdev}$, $\val{AxisWdev}$ and $\val{HighWdev}$. The
    step-doubling estimates are $\val{LowErr}$, $\val{AxisErr}$ and
    $\val{HighErr}$, and the largest of the flights is $\val{errStep}$. The
    panel-size estimate is $\val{errMesh}$. The sampled clearance is
    $\val{refDmin}$, outside the shell $\val{contact}$ and outside the charge
    radius $\val{Rcharge}$. The image-force ratio reaches $\val{refImage}$.
  \item\label{st:rad} \emph{The radiation.} Simpson's rule on steps of $h/2$
    gives $\val{LowRad}$, $\val{AxisRad}$ and $\val{HighRad}$. The largest
    fraction of $T_0$ is $\val{radFracMax}$. Halving the quadrature step again
    changes the largest of those integrals by at most the relative amount
    $\val{LowRadErr}$. These integrals are not rows of the comparison.
\end{procedure}

\begin{result}[The reference flights]
All three ions enter the detector: the downward ray through the bottom edge
at $(\val{LowX}, \val{LowY})$, the on-axis ray through the left edge at
$(\val{AxisX}, \val{AxisY})$, and the upward ray through the top edge at
$(\val{HighX}, \val{HighY})$. The level is solved.
\end{result}

\section{Comparison with the game}\label{sec:comparison}

\comparisontable

Twenty-four comparisons are in the table: outcome, time and end position for
each of the three reference rays, and the same four numbers with the charges
omitted. The tolerance on a state is $\val{tolFlight}$. Ten times the
panel-size estimate $\val{errMesh}$ lies above $10^{-9}$, and it is larger
than the step-doubling estimate $\val{errStep}$, so the tolerance is that
product. The tolerance on a state was fixed from these estimates before the
comparison.

The largest absolute difference is the lower reference abscissa,
$\val{cmpLowXDiff}$. The lower reference time differs by $\val{cmpLowTDiff}$.
Both lie under the panel-size estimate $\val{errMesh}$, and so under the
tolerance. That estimate is large because an end state moves by
$\val{errMesh}$ when the panels double, and the lower ray meets the detector
along its bottom edge, where the meeting point slides. The step estimate
$\val{errStep}$ is far below those differences: the panels, not the time
step, set them. The outcomes agree.

The lower reference ordinate agrees exactly, and it is the detector's bottom
edge. The on-axis reference abscissa and the bare on-axis abscissa agree
exactly, and each is the detector's left edge. The bare upper abscissa agrees
exactly, and it is the world's right edge. The bare lower abscissa differs by
$\val{cmpOffLowXDiff}$, under the step estimate. The upward reference ordinate
differs by $\val{cmpHighYDiff}$: our march meets the top edge, and the game's
meeting point is just inside that edge.

\section{Other solutions}\label{sec:others}

The empty placement in \code{placements.json} omits both charges and does not
solve the level. The lens alone remains. The on-axis ray still meets the
\val{OffAxisFace} face of the detector at $(\val{OffAxisX}, \val{OffAxisY})$,
at $t = \val{OffAxisT}$. The downward ray leaves through the \val{OffLowFace}
edge of the world at $(\val{OffLowX}, \val{OffLowY})$, at $t = \val{OffLowT}$,
at an angle $\val{OffLowAng}$ degrees from $+x$. The ordinate lies below
$\val{yDetLo}$. The upward ray leaves through the \val{OffHighFace} edge at
$(\val{OffHighX}, \val{OffHighY})$, at $t = \val{OffHighT}$, at an angle
$\val{OffHighAng}$ degrees from $+x$, above $\val{yDetHi}$.

The sampled clearance with the charges omitted is $\val{offDmin}$, and the
image-force ratio reaches $\val{offImage}$. The two reference charges are what
brings the angled rays from those two misses onto the detector. This document
integrates the reference and the empty placement.

\section{What the game computes}\label{sec:engine}

The middle electrodes are a fixed potential and the outer electrodes are
grounded (\physref{2.7}). The surface charge is the engine's triangle
collocation, the same boundary condition as
\xeqref{deflection_plates}{eq:collocation}, with the placed charges on the
right-hand side. This document uses the equal rectangles of the deflection
plates. The end states move by $\val{errMesh}$ when that edge doubles, and
the lower reference abscissa in \S\ref{sec:comparison} lies under that
estimate.

The ion's image force is omitted on both sides. The bound sampled on the
reference flights reaches $\val{refImage}$, under the engine's limit
$10^{-10}$. The flight in the engine is the adaptive method of \physref{5}
and \physref{6}, verified as in \physref{7}. Our step estimate is
$\val{errStep}$. Radiation is integrated and not applied to the momentum
(\xsecref{foundations}{sec:radiation}, \physref{8}).

\section{Further reading}\label{sec:reading}

The rectangle integral, the collocation and the flight energy are derived for
the deflection plates (\xeqref{deflection_plates}{eq:rectangle},
\xeqref{deflection_plates}{eq:collocation},
\xeqref{deflection_plates}{eq:flight-energy}). The symmetric aperture lens is
derived for Tune the lens (\xeqref{tune_the_lens}{eq:axis},
\xeqref{tune_the_lens}{eq:slow}). The motion and the Runge--Kutta step are in
the Foundations (\xsecref{foundations}{sec:numerics}). The engine's
electrodes are \physref{2.7}.

\levelfileappendix

\end{document}
